\documentclass[aps, prx, nofootinbib, preprintnumbers, superscriptaddress]{revtex4-2}

\usepackage[utf8]{inputenc}
\usepackage{graphicx}
\usepackage{amsmath, amssymb, amsfonts, amsthm, bm, nicefrac}
\usepackage{multirow} 
\usepackage{tabularx} 
\usepackage{array}
\usepackage{units}
\usepackage{tensor} 
\usepackage{braket}
\usepackage{hyperref}
\usepackage{dcolumn}
\usepackage{bm}
\usepackage{enumitem}
\setlist[itemize,2]{label=\textopenbullet} % itemize의 두 번째 레벨의 기호를 \textbullet로 설정
\usepackage{booktabs} % 테이블 선을 위한 패키지 추가


\usepackage{xcolor}
\usepackage{soul}
\usepackage[export]{adjustbox}
\usepackage{tikz}
\usetikzlibrary{positioning}



% Define global settings and functions before \begin{document}
\definecolor{SkyBlue}{rgb}{0.5, 0.8, 1} 
\definecolor{MyPink}{rgb}{1, 0.682, 0.843}
\definecolor{customblue}{HTML}{1f77b4}
\definecolor{customorange}{HTML}{ff7f0e}
\definecolor{customgreen}{HTML}{2ca02c}
\definecolor{custompurple}{HTML}{C4A3F0}





% Define a custom function for magnetic ordering globally
\newcommand{\magneticOrdering}[5]{
    % Arguments: { #1: grid row }{ #2: grid column }{ #3: angle theta in degrees }{ #4: length of the arrow }{ #5: color of the arrow }

    % Calculate the grid point's actual position
    \pgfmathsetmacro\xa{#2 * \size - #1 * 0.5 * \size}
    \pgfmathsetmacro\ya{#1 * \size * 0.866} % Height of equilateral triangle

    % Calculate arrow end point
    \pgfmathsetmacro\dx{#4 * cos(#3)} % x-component of the arrow
    \pgfmathsetmacro\dy{#4 * sin(#3)} % y-component of the arrow

    % Draw arrow
    \draw[->, line width=1.5pt] (\xa, \ya) -- ++(\dx, \dy);

    % Draw red dot
    \fill[#5] (\xa, \ya) circle(0.2);
}



\setcounter{secnumdepth}{2}
\makeatletter 
\renewcommand\twocolumngrid{
	\def\footnoterule{
		\dimen@\skip\footins\divide\dimen@\thr@@
		\kern-\dimen@\hrule width.5in\kern\dimen@}
	\do@columngrid{mlt}{\tw@}
 }

\makeatother  
\newtheorem{theorem}{Theorem}[section]
\newtheorem{lemma}[theorem]{Lemma}

\newtheorem{definition}{Definition}[section]
\newtheorem{prop}{Proposition}
\newtheorem{axiom}{Axiom}
\newtheorem{claim}{Claim}
\newtheorem{corollary}{Corollary}[theorem]

\hypersetup{
    colorlinks=true,
    linkcolor=blue,
    filecolor=magenta,      
    urlcolor=blue,                                         %% 
    pdftitle={Stability of Chirality for mixed states}, %% PDF의 제목
    pdfpagemode=FullScreen,
    }

\newcommand{\SP}[1]{{\color{cyan}\footnotesize{(SP) #1}}}
\newcommand{\NBCPatom}{$\mathrm{Na_2 Ba Co (PO_4)_2}$}
\newcommand{\NBCP}{{\sf NBCP}}

\newcommand{\kvec}{\mathbf{k}}
\newcommand{\bmdelta}{\bm{\delta}}

\newcommand{\ha}{\hat{a}}
\newcommand{\hc}{\hat{c}}
\newcommand{\hI}{\hat{\mathbb{I}}}
\newcommand{\hP}{\hat{P}}

\newcommand{\rD}{\mathrm{D}}
\newcommand{\rU}{\mathrm{U}}
\newcommand{\ri}{\mathrm{i}}


\newcommand{\bfr}{\mathbf{r}}




\newcommand{\cA}{\mathcal{A}}
\newcommand{\cB}{\mathcal{B}}
\newcommand{\cC}{\mathcal{C}}
\newcommand{\cD}{\mathcal{D}}
\newcommand{\cE}{\mathcal{E}}
\newcommand{\cG}{\mathcal{G}}
\newcommand{\cH}{\mathcal{H}}
\newcommand{\cJ}{\mathcal{J}}
\newcommand{\cN}{\mathcal{N}}
\newcommand{\cM}{\mathcal{M}}
\newcommand{\cP}{\mathcal{P}}
\newcommand{\cR}{\mathcal{R}}
\newcommand{\cQ}{\mathcal{Q}}
\newcommand{\cU}{\mathcal{U}}


\newcommand{\Id}{\mathrm{Id}}

\newcommand{\bth}{\bm{\theta}}


\newcommand{\Tr}{\mathrm{Tr}}
\newcommand{\Pf}{\mathrm{Pf}}


\newcommand{\Theme}[1]{{\vspace{0.5cm}\noindent\textbf{#1}}}

\begin{document}

\title{Note: Linear Spin Wave Theory}
\author{Sung-Min Park}
\begin{abstract}
    In this article, we summarize the linear spin wave theory (LSWT) described in the literature.
\end{abstract}
\maketitle

\tableofcontents

\newpage
\section*{Summary of Notation and Symbols}

We summarize the symbols and notation used in this notes.

\begin{table*}[!h]
    \centering
    \caption{Summary of Notation and Symbols in Linear Spin Wave Theory}
    \begin{tabular}{c|p{14cm}}
    \hline\hline
    \multicolumn{2}{c}{\bf Spin Hamiltonian and Related Quantities} \\
    \hline
    $\hat{H}$ & Spin Hamiltonian \\
    $J_{ij}^{\alpha\beta}$ & Exchange coupling between sites $i$ and $j$ for spin components $\alpha$ and $\beta$ \\
    $h_j^{\alpha}$ & External magnetic field at site $j$ along direction $\alpha$ \\
    $E_{\rm cl}$ & Classical ground state energy \\
    $\mathbf{S}_j$ & Classical spin vector at site $j$ \\
    $\hat{\mathbf{S}}_j$ & Quantum spin operator at site $j$ \\
    $\delta\hat{\mathbf{S}}_j$ & Quantum fluctuation operator ($\hat{\mathbf{S}}_j = \mathbf{S}_j + \delta\hat{\mathbf{S}}_j$) \\
    $\widetilde{\mathbf{S}}_j$ & Spin operator in the local reference frame \\
    \hline
    \multicolumn{2}{c}{\bf Index Conventions and Coordinate Systems} \\
    \hline
    $i, j$ & Unit cell indices \\
    $\mu, \nu, \sigma, \rho$ & Sublattice indices (typically $a, b, c, \ldots$) \\
    $I = (i, \mu), J = (j, \nu)$ & Composite site indices combining unit cell and sublattice \\
    $\mathbf{r}_i$ & Position of unit cell $i$ \\
    $\boldsymbol{\delta}_\mu$ & Position of sublattice $\mu$ within the unit cell \\
    $\mathbf{R}_I = \mathbf{r}_i + \boldsymbol{\delta}_\mu$ & Position of spin at site $I = (i, \mu)$ \\
    $\alpha, \beta$ & Spin component indices ($x, y, z$ or $+, -, 0$) \\
    $\mathbf{R}_j$ & Rotation matrix from laboratory frame to local frame at site $j$ \\
    $\hat{\mathbf{e}}_j^{\alpha}$ & Basis vectors of local frame at site $j$ ($\alpha = +, -, 0$) \\
    \hline
    \multicolumn{2}{c}{\bf Boson Operators and Transformations} \\
    \hline
    $\hat{a}_j, \hat{a}_j^{\dagger}$ & Bosonic annihilation and creation operators \\
    $\hat{n}_j = \hat{a}_j^{\dagger}\hat{a}_j$ & Boson number operator \\
    $\hat{b}_{\mathbf{k}\mu}, \hat{b}_{\mathbf{k}\mu}^{\dagger}$ & Momentum-space bosonic operators for sublattice $\mu$ \\
    $\hat{\beta}_{\mathbf{k}\mu}, \hat{\beta}_{\mathbf{k}\mu}^{\dagger}$ & Diagonalized magnon operators (after Bogoliubov transformation) \\
    $\Psi_{\mathbf{k}} = (\hat{\mathbf{b}}_{\mathbf{k}}, \hat{\mathbf{b}}_{-\mathbf{k}}^{\dagger})^T$ & Nambu spinor in momentum space \\
    $\widetilde{\Psi}_{\mathbf{k}} = (\hat{\boldsymbol{\beta}}_{\mathbf{k}}, \hat{\boldsymbol{\beta}}_{-\mathbf{k}}^{\dagger})^T$ & Diagonalized Nambu spinor \\
    \hline
    \multicolumn{2}{c}{\bf Matrices and Transformation Operators} \\
    \hline
    $\mathsf{H}_{\mathbf{k}}$ & Hamiltonian matrix in BdG form \\
    $\mathsf{A}_{\mathbf{k}}, \mathsf{B}_{\mathbf{k}}$ & Blocks of the Hamiltonian matrix $\mathsf{H}_{\mathbf{k}}$ \\
    $\mathsf{T}_{\mathbf{k}}$ & Paraunitary transformation matrix for Bogoliubov diagonalization \\
    $\mathsf{E}_{\mathbf{k}}$ & Diagonal matrix of magnon energies \\
    $\mathsf{N}_{\mathbf{k}}(t)$ & Time-dependent correlation matrix in magnon basis \\
    $\mathsf{U}^{\alpha}$ & Spin projection operator in laboratory frame \\
    $\mathsf{S}_{\mathbf{k}}$ & Matrix encoding spin magnitudes and sublattice phase factors \\
    $\mathsf{C}^{\alpha\beta}(\mathbf{k}, t)$ & Correlation matrix in laboratory frame \\
    $\mathsf{J}$ & Symplectic metric tensor \\
    \hline
    \multicolumn{2}{c}{\bf Physical Quantities} \\
    \hline
    $E_{\mathbf{k},\mu}$ & Magnon energy for band $\mu$ at momentum $\mathbf{k}$ \\
    $n_{\mathbf{k},\mu} = 1/(e^{\beta E_{\mathbf{k},\mu}}-1)$ & Bose-Einstein distribution function \\
    $E_0$ & Zero-point energy \\
    $\mathcal{C}^{\alpha\beta}_{\mu\nu}(\mathbf{k}, t)$ & Sublattice spin-spin correlation function \\
    $\mathcal{S}^{\alpha\beta}(\mathbf{k}, \omega)$ & Dynamic structure factor \\
    $\mathcal{A}^{\alpha\beta}(\mathbf{k}, \omega)$ & Spectral function \\
    $\Omega_{n,\mathbf{k}}$ & Berry curvature of the $n$-th magnon band \\
    $C_n$ & Chern number of the $n$-th magnon band \\
    \hline\hline
    \end{tabular}
\end{table*}

\newpage
